\section{Literature Review}\label{sec:lit}

Prepayment theory begins with ``ruthless exercise''---instantaneous refinancing as conditions dictate \citep{dunn1981}. It then admits frictions. Within a valuation framework, \citet{schwartz1989} fit an empirical prepayment function to aggregate GNMA experience by maximum-likelihood hazard, the econometric ancestor of my loan-level survival approach (Section~\ref{sec:hazard}). \citet{richard1989} formalize industry practice: prepayment as an empirical function of refinancing incentive, seasoning, seasonality, and a pool's own prepayment history. \citet{stanton1995} add heterogeneous transaction costs, which smooth the refinancing cliff into a predictable S-curve. \citet{agarwal2013} sharpen the rational benchmark into a closed-form optimal refinancing threshold that scores observed behavior. \citet{chernov2018} recover from MBS market prices a function driven by macroeconomic turnover alongside the rate response. My rate-inelastic null operationalizes that turnover-versus-rate-response split on the cash-flow side. \citet{quigley1987,quigley2002} established lock-in as a persistent friction on mobility in rising-rate environments. \citet{ferreira2010} gave early large-sample evidence that rising mortgage rates and negative home equity each depress owner mobility. The latter is an equity-lock-in channel, distinct from and complementary to the rate channel this paper models. The equity dimension is left unmodeled here, a boundary disclosed in Section~\ref{sec:pathb}.

The 2022--2025 rate shock generated a large empirical literature documenting the breakdown of standard prepayment assumptions. \citet{fonseca2024} documented the steep decline in U.S. mobility and refinancing. \citet{fonseca2026} extend that evidence to general equilibrium, where locked-in owners who do not downsize raise net housing demand, offsetting part of the price decline a rate rise would otherwise produce. This paper neither models nor prices that channel. \citet{batzer2024} estimate roughly \$2.4 trillion of below-market financing value held by U.S. fixed-rate borrowers---retained by staying, forfeited on relocation---and roughly 1.33 million sales prevented over 2022Q2--2023Q4 (1.72 million through 2024Q2 in their revision). \citet{liebersohn2024} put it in hazard form: each percentage point of prevailing rate above a borrower's locked-in coupon reduces quarterly mobility probability by 5.5\% to 7.7\% depending on specification. My Path B microsimulation adopts that elasticity band as calibration. \citet{graybill2026}, a Federal Reserve Bank of Philadelphia working paper, independently estimate that margin from transaction records in a housing-market-equilibrium setting. They find deep rate gaps cut the sale hazard by roughly 7--11\% per percentage point. That setting is where lock-in's most-discussed housing consequence sits: a household that does not list withholds a unit. The incidence falls on first-time buyers and renters, not the locked-in owner, a distributional consequence this design measures nowhere. The finding corroborates the band from data that never touch credit records. But it sits at or above the band, not at its centre: only the 7\% low end falls inside 5.5\% to 7.7\%, and its midpoint lies above the high edge (Section~\ref{sec:pathb}).

A related strand asks how much of the mobility decline a rational rate-gap framework explains, leaving the rest an unexplained residual. \citet{aladangady2024} hold the 2021 rate environment fixed in a counterfactual rate-gap distribution and correct for selection bias from forgone refinancing; they find rate-gap-driven lock-in explains only 44\% of the actual decline in mortgage borrower mobility between 2021 and 2022. Section~\ref{sec:identification} maps that 44\% share into this design's own units as a calibration target for the scaled-null variant (run \texttt{scaled\_null\_housing\_activity}). My decomposition in Sections~\ref{sec:abm} and~\ref{sec:hazard} is methodologically analogous, but it runs on dynamically modeled CPR and cash-flow shortfalls rather than mobility rates, with two structurally distinct estimators. \citet{gerardi2024} similarly score a structural search-and-matching lock-in model against empirical lifecycle migration, finding it accounts for most, but not all, of the observed disruption to housing market liquidity.

A separate behavioral finance literature grounds what such residuals likely represent. \citet{keys2016} find more than 20\% of U.S. borrowers fail to refinance when it is unambiguously in their financial interest, attributing this to procrastination and uneven financial sophistication. \citet{defusco2020} add a structural counterpart: employment-documentation requirements and out-of-pocket closing costs bind hardest in recessions and inhibit refinancing pass-through. On Danish administrative data, \citet{andersen2020} find inattention and inertia, not transaction costs, behind much failure to act on refinancing opportunities. \citet{andersen2022}, on Danish housing data, show reference dependence and loss aversion relative to the original purchase price, with losses weighted roughly 2.5 times more heavily than equivalent gains. Across administrative, survey, and field-experimental data, \citet{johnson2019} find that of the behavioral factors tested, suspicion of lenders' motives---not standard rational frictions such as time preference or risk aversion---is most consistently related to households' failure to accept profitable refinancing offers, which further supports a behavioral reading of the residual. \citet{bissiri2014} formalize the distinction, defining ``behavioral risk'' as the share of prepayment decisions not driven by a fully rational financial criterion. \citet{perotti2024} applies that operationalization more recently, modeling behavioral uncertainty as a stochastic overlay on an otherwise rational prepayment framework and deriving robust hedging strategies that bound, not eliminate, the exposure it creates.

A further strand treats prepayment as a loan-level survival process shaped by pool composition and servicer behavior, not by borrower optimization. Prepayment ``burnout'' is a standard feature of industry MBS valuation models and one of the four factors of \citepos{richard1989} canonical prepayment function, alongside the Public Securities Association seasoning convention. Burnout is the tendency for a pool's aggregate refinancing sensitivity to decline over time as its most rate-sensitive borrowers prepay earlier and self-select out of the remaining pool. Yet burnout is difficult to reconcile with a fixed, representative-agent population. Recent theory formalizes why. In a heterogeneous pool the observed pool-level hazard is the survival-weighted mean of the individual hazards. Its drift carries a negative selection term proportional to the cross-sectional dispersion of those hazards \citep[arXiv preprint]{lesniewski2026}, so pool speeds decline with seasoning even when no individual hazard changes. My hazard framework (Section~\ref{sec:hazard}) is designed to accommodate this loan-level survival structure.

A separate fixed-income literature highlights the systemic risks the lock-in effect generates. \citet{jiang2023} documented that duration extension in bank asset portfolios, prominently including low-coupon MBS, generated roughly \$2 trillion in unrecognized mark-to-market losses across the U.S. banking sector, contributing to self-fulfilling solvency runs such as the collapse of Silicon Valley Bank. \citet{smith2023} demonstrated that QT relies heavily on predictable MBS runoff schedules, which these duration extensions structurally break.

Closest to this paper, the Federal Reserve's own staff and officials have already described the mechanism and measured its first-order footprint. \citet{na2024} report that ``a vast majority of the agency MBS held in the SOMA are far out-of-the-money to refinance and may be subject to the `lock-in' effect.'' They further report that because ``monthly principal payments have been below the \$35 billion-per-month redemption cap \ldots the size of the SOMA agency MBS portfolio has decreased at a much slower pace than in the case of a binding cap.'' They put June 2022 to June 2024 runoff at roughly \$450 billion against the \$820 billion a binding cap would have retired, an implied shortfall near \$370 billion. The same note reports the composition of those below-cap payments: scheduled principal accounts for roughly half, with prepayments mostly due to turnover making up the remainder. That is an external, staff-measured anchor for the scheduled-plus-turnover decomposition my $\beta_1 = 0$ null quantifies. The anchor and my null have different denominators. Na et al.'s composition describes realized principal \emph{payments}: roughly half scheduled, half turnover-driven prepayment. My null is built from the same two ingredients, scheduled amortization plus a turnover floor calibrated from realized, partly behavioral turnover. My 88.7\% (in-sample; 85.7\% at the off-window floor) instead describes the null's share of the shortfall \emph{against the caps}. A payments mix of scheduled-plus-turnover that runs below \$35 billion in essentially every month is the state in which that shortfall accrues mechanically. Their composition and my null recovery are therefore the same fact read against two denominators, and read on the composition the null is built to reproduce. \citet{perli2024}, then Manager of the System Open Market Account, stated that the cap ``has not been binding \ldots resulting in very slow prepayments,'' and \citet{hammack2025} observed that some of the low-coupon holdings ``could be very long lived.'' New York Fed staff had projected the below-cap regime from its first month: a September 2022 staff post put monthly agency MBS principal payments at \$20 to \$30 billion over the following months, below the \$35 billion cap that took effect that month \citep{nyfed2022}. I treat the qualitative mechanism as established by these sources. Four features separate my exercise from theirs. First, this characterization attributes the below-cap runoff to lock-in without isolating what the elasticity contributes, as distinct from the scheduled amortization and baseline turnover that would have undershot the cap under any rate-driven prepayment response above that baseline. Sections~\ref{sec:abm} and~\ref{sec:hazard} supply that decomposition. They find the identified lock-in margin bounded between $+2.3$ and $+9.1$ points of the benchmark at a turnover floor measured off the locked-in population. That bound is a wild-cluster interval on the floor read's 31 clusters; the narrower percentile read under-covers. At the mid-grid anchor inside that interval the margin is $+5.6$ percentage points ($+9.2$ at the in-window production calibration), and Section~\ref{sec:identification} qualifies that margin. Second, I extend the same shortfall-against-cap construction from the \$450 billion-through-June-2024 quantity to the full June 2022 to November 2025 window, where it reaches \$764.7 billion (Section~\ref{sec:method-benchmark}). Third, I measure the shortfall against the Federal Reserve's own ex-ante projection as well as against the caps, separating the component the Federal Reserve itself anticipated from the genuine surprise (Section~\ref{sec:method-benchmark}). Fourth, the decomposition turns the staff observation into a statement about cap design: scheduled amortization and a turnover floor read from realized, partly behavioral turnover undershoot a \$35 billion monthly ceiling in essentially every month, so a cap at that level could not have bound under any rate response above that floor (Section~\ref{sec:identification}). A more recent \citet{eyal2026} FEDS Note also decomposes the Federal Reserve's balance-sheet reduction, but into macroeconomic drivers (inflation, real GDP growth, and active SOMA reductions relative to nominal GDP), not into the loan-level sources of the MBS runoff shortfall itself, so it does not speak to the lock-in-versus-baseline question I take up here.

The Danish mortgage securitization system offers a tested structural alternative to negative convexity. \citet{berg2018} set out its ``Balance Principle'' of strict one-to-one match funding; under those rules \citet[\S3.3.2]{berger2026}, in a July 2026 SSRN working draft, estimate household mobility close to flat in the coupon gap, because borrowers may repurchase their outstanding mortgages at market value, neutralizing extension risk for the bondholder. The mechanics are institutional, not notional, and fix what a transplant must replicate: beyond penalty-free prepayment at par, a Danish mortgagor may retire the loan by buying the corresponding bonds in the secondary market and \emph{delivering} them to the mortgage bank (an option the U.S. contract does not carry), and the Danish contract has no due-on-sale clause \citep{frankel2004}. That is an open-market purchase settled by delivery, not a lender-side discount: its cost to the borrower is the bond's market price, its availability depends on a liquid secondary market in the specific series, and the bondholder is a \emph{voluntary seller} at market price, not called at par. Transplanted onto a passthrough holder the analogue is an outright sale at a discount, so a par-denominated redemption cap is not the natural scorer for the Danish leg, hence bracketing it under face and cash accounting. \citet{svenstrup2006} set out the design's wider reform implications. The flatness is a slope result on Danish data, so Section~\ref{sec:abm-danish} is careful what it licenses: it imports the slope but anchors the level at the U.S. hazard, Berger et al.\ estimating no moving level for U.S. households under a Danish payoff rule. The strand supplies rule and slope, not a priced U.S. counterfactual; Section~\ref{sec:abm-danish} scores the rule-only transplant on the SOMA book directly, an order of magnitude below a mechanism-extrapolated read, and Appendix~\ref{app:params} prices the buyback credit's incidence, which decides whether the institutional gap survives the accounting.

A final strand situates lock-in within the theory of monetary policy transmission. \citet{berger2021} show the Federal Reserve's power to stimulate through rate cuts is path-dependent, since prepayment depends on the outstanding distribution of in-the-money refinancing incentives. I extend that to tightening. The path dependence that weakens easing after a low-rate period weakens tightening after one too: the resulting stock of deeply in-the-money, rate-locked mortgages cannot prepay, however aggressively the central bank then raises rates. \citet{eichenbaum2022} formalize this state dependence for the easing direction, where policy potency depends on the distribution of outstanding rate gaps. \citet{beraja2019} add regional evidence that the refinancing channel's strength varies systematically with the local distribution of in-the-money mortgages. \citet{dimaggio2020} give direct evidence for the easing-direction channel this paper mirrors: the Federal Reserve's QE1 agency-MBS purchases transmitted to the real economy specifically by inducing refinancing. The QT setting also links to how balance-sheet operations transmit through the supply of long-duration assets. There \citet{krishnamurthy2011} document several distinct quantitative-easing channels, one specific to agency MBS and running through prepayment. \citet{vayanos2021} formalize, under preferred-habitat demand, why shifts in the outstanding maturity and supply of held bonds move the term structure at all. Lock-in is thus an input to that balance-sheet-supply channel, not separate from it. \citet{du2024} assemble the cross-country QT record---seven central banks, active and passive runoff---finding announcement effects concentrated in yields. The MBS runoff shortfall decomposed here is that record's passive-runoff margin, measured at the portfolio cash-flow level. \citet{vickery2013} trace the TBA market's liquidity premium specifically to the homogeneity of deliverable pools, with TBA-eligible securities trading at significantly lower cost than otherwise similar specified pools. That homogeneity requirement is the institutional constraint that, in Section~\ref{sec:conclusion}, I argue makes the Danish market-value repurchase structure hard to reconcile with the current U.S. secondary market absent a broader redesign of agency securitization rules.

My two estimators sit at the transparent end of two active methodological literatures. On the survival side, prepayment and default hazards are increasingly fit nonparametrically or by machine learning: random survival forests \citep{ishwaran2008}, competing-risks Bayesian additive regression trees \citep{sparapani2020}, and deep neural-network mortgage hazard models \citep{sadhwani2021}. The last benchmarks what flexibility buys at scale. On origination and monthly performance records for over 120 million U.S. mortgages, the deep network of \citet{sadhwani2021} learns highly nonlinear, interaction-rich prepayment and default responses without hand-specified functional form. Those responses include the burnout-type history dependence Path~A carries as an explicit stratum-level regressor and the ABM probes by pre-committed falsification (Sections~\ref{sec:patha} and~\ref{sec:abm-results}). I instead estimate a fixed-effects Poisson hazard and a literature-calibrated microsimulation, with functional forms and calibration inputs written out in full. What I need is not maximal predictive fit but an auditable decomposition. The null of Section~\ref{sec:identification} that switches the lock-in elasticity off is interpretable only if every term it holds fixed is visible, and a black-box hazard would obscure those terms. Those machine-learning survival methods remain a natural robustness avenue for the loan-level fit, and Section~\ref{sec:patha} probes them at their tabular end. Gradient-boosted Poisson and isotonic-recalibration comparators are each given Path A's information set on the identical pre-committed holdout, and each recalibrates levels without gaining cross-cell predictive power. Deep loan-level variants remain future work on data-availability grounds. A complementary strand attacks nonstationarity, not nonlinearity. Landmarking re-anchors a simple hazard among survivors at successive prediction times \citep{vanhouwelingen2007,putter2022}. Recent credit-risk work pairs landmark targets with explicit probability recalibration to hold survival predictions stable under covariate and base-rate drift, on the same Freddie Mac performance data my hazard paths use \citep{peng2026}. My design meets that concern with fixed, visible machinery: stratum fixed effects, calendar-month dummies, the dynamic friction index of Section~\ref{sec:patha}. It prices the residual time-variation it cannot remove by the temporal moving-block bootstrap reported there. On the agent-based side my calibration is fixed by the pre-committed parameterization of Appendix~\ref{sec:method-abm}, not by an automated estimator. It sits alongside a maturing literature on data-driven ABM calibration and estimation: simulated-minimum-distance and Bayesian methods \citep{grazzini2017,platt2020}, the empirical-validation taxonomy of \citet{fagiolo2007}. Within that literature the pre-committed stress-test design of Section~\ref{sec:abm} is one validation strategy among several. The paired design-crossing of Section~\ref{sec:robustness-crossdesign} is offered as a validation design in its own right: the ABM's decision rule re-run on real loan covariates, the same survival structure run on a fully synthetic population, interpretive thresholds fixed before the cross-design runs. Its nearest relative in that literature is indirect inference, which minimizes the distance between an auxiliary model's parameters fitted on real and on simulated data \citep{platt2020}. The pair here crosses designs not to estimate parameters but to attribute a result gap to modeling paradigm versus data source. It is not, to my knowledge, a named member of the validation taxonomy of \citet{fagiolo2007} or of the estimation approaches surveyed by \citet{grazzini2017} and \citet{platt2020}.